\section{Elliptic normalization and arithmetic structure}\label{app:elliptic}

This appendix develops the arithmetic--geometric structure underlying the quartic spectral curve $\Pi(\lambda,y)=0$.
We show that $\Pi$ defines an elliptic curve of conductor~$37$, that the branch locus of Theorem~\ref{thm:disc-factor} is the ramification image of a degree-$4$ morphism on this curve, and we record the Galois-theoretic, divisor-theoretic, and local-analytic consequences that follow from classical algebraic geometry and number theory.

\subsection{Ellipticization of the quartic spectral curve}

Let $C\subset \mathbb{A}^2_{\lambda,y}$ be the affine curve cut out by $\Pi(\lambda,y)=0$.

\begin{theorem}[Ellipticization]\label{thm:ellipticization}
Under the polynomial change of coordinates
\[
(\lambda,y)\longmapsto (\lambda,\xi:=y-\lambda^2),
\]
the affine curve $C$ is carried isomorphically onto the affine Weierstrass curve
\begin{equation}\label{eq:elliptic-E}
E:\qquad \xi^2+\xi=\lambda^3-\lambda.
\end{equation}
Consequently, the smooth projective model of $C$ is an elliptic curve over $\QQ$.
Moreover, in the integral minimal model \eqref{eq:elliptic-E} one has
\[
\Delta(E)=37,\qquad j(E)=\frac{110592}{37},
\]
and $E$ lies in the conductor-$37$ isogeny class (Cremona label $37\mathrm{a}1$).
\end{theorem}

\begin{proof}
Viewing $\Pi(\lambda,y)$ as a quadratic in $y$ gives
\[
y^2+(1-2\lambda^2)y+\bigl(\lambda^4-\lambda^3-\lambda^2+\lambda\bigr)=0.
\]
Substituting $y=\lambda^2+\xi$ yields
\[
(\xi+\lambda^2)^2+(1-2\lambda^2)(\xi+\lambda^2)+(\lambda^4-\lambda^3-\lambda^2+\lambda)=0.
\]
Expanding: $\xi^2+2\lambda^2\xi+\lambda^4+\xi+\lambda^2-2\lambda^2\xi-2\lambda^4+\lambda^4-\lambda^3-\lambda^2+\lambda=0$, which simplifies to $\xi^2+\xi=\lambda^3-\lambda$.
Since the coordinate change is invertible with inverse $(\lambda,\xi)\mapsto(\lambda,y=\lambda^2+\xi)$,
it restricts to an isomorphism of affine curves.

For the invariants, \eqref{eq:elliptic-E} is a Weierstrass equation with coefficients
$(a_1,a_2,a_3,a_4,a_6)=(0,0,1,-1,0)$.
By the standard discriminant formula~\cite{Silverman2009}, $\Delta=37$ and $j=110592/37$.
The curve is minimal at every prime (since $\Delta=37$ is squarefree), and the conductor equals $37$ since $E$ has multiplicative reduction at $p=37$ and good reduction elsewhere.
\end{proof}

\subsection{The parameter \texorpdfstring{$y$}{y} as a degree-\texorpdfstring{$4$}{4} map and its ramification}

On $E$ define the rational function
\begin{equation}\label{eq:y-as-function}
y=\lambda^2+\xi\in \QQ(E),
\end{equation}
and let $\varphi:E\to \mathbb{P}^1_y$ be the associated morphism.

\begin{theorem}[Ramification structure]\label{thm:ramification-image}
The map $\varphi:E\to \mathbb{P}^1_y$ has degree $4$.
Its ramification data are as follows:
\begin{enumerate}
\item[(i)] Over $y=\infty$: a single totally ramified point $O\in E$ with $e_O=4$.
\item[(ii)] Over $y=0$: one point with $e=2$ (at $\lambda=1$, $\xi=-1$) and two unramified points.
\item[(iii)] Over $y=1$: one point with $e=2$ (at $\lambda=-1$, $\xi=0$) and two unramified points.
\item[(iv)] Over each root $y_\star$ of $c(y)=0$: one point with $e=2$ and two unramified points.
\end{enumerate}
The finite branch locus is exactly $\{y=0\}\cup\{y=1\}\cup\{c(y)=0\}$, in agreement with Theorem~\ref{thm:disc-factor}.
The total ramification satisfies the Riemann--Hurwitz formula:
\[
2g(E)-2=4\cdot(2\cdot 0-2)+\sum(e_P-1)=-8+(4-1)+(2-1)\cdot 5=0,
\]
confirming $g(E)=1$.
\end{theorem}

\begin{proof}
At the point at infinity $O$ of the Weierstrass model \eqref{eq:elliptic-E}, the standard pole orders are $\operatorname{ord}_O(\lambda)=-2$ and $\operatorname{ord}_O(\xi)=-3$.
Therefore $\operatorname{ord}_O(y)=\operatorname{ord}_O(\lambda^2+\xi)=\min(-4,-3)=-4$,
so $y$ has a unique pole of order $4$, giving $\deg(\varphi)=4$ with $y=\infty$ totally ramified ($e_O=4$, contributing $e_O-1=3$).

For finite ramification, differentiate \eqref{eq:elliptic-E} to get $(2\xi+1)\,\dd\xi=(3\lambda^2-1)\,\dd\lambda$.
From $y=\lambda^2+\xi$:
\[
\dd y=\frac{2\lambda(2\xi+1)+3\lambda^2-1}{2\xi+1}\,\dd\lambda.
\]
Thus $\dd y=0$ iff $2\lambda(2\xi+1)+3\lambda^2-1=0$, i.e., $\xi=(1-2\lambda-3\lambda^2)/(4\lambda)$ for $\lambda\neq 0$.
Substituting into \eqref{eq:elliptic-E} and clearing denominators yields
\[
(\lambda-1)(\lambda+1)(16\lambda^3-9\lambda^2+1)=0.
\]
This gives five finite ramification points (each simple, contributing $e-1=1$), mapping to branch values:
$\lambda=1\mapsto y=0$; $\lambda=-1\mapsto y=1$; three roots of $16\lambda^3-9\lambda^2+1\mapsto$ three roots of $c(y)$.
Total ramification: $3+5\cdot 1=8=-8+\deg\cdot(2\cdot 0)$? Let us verify: $2g(E)-2=0$; $\deg\varphi(2g(\mathbb{P}^1)-2)=4\cdot(-2)=-8$; $\sum(e_P-1)=3+5=8$. Indeed $0=-8+8$. \checkmark
\end{proof}

\begin{remark}[Hurwitz passport]\label{rmk:hurwitz}
The ramification data of $\varphi$ can be recorded as a \emph{Hurwitz passport}---the tuple of cycle types of the monodromy permutations at each branch value.
With $\deg\varphi=4$ and the ramification computed above, the passport over the five finite branch values is:
$y=0$: $(2,1,1)$; $y=1$: $(2,1,1)$; each root of $c(y)$: $(2,1,1)$.
Over $y=\infty$: $(4)$.
The monodromy group generated by these permutations is $S_4$ (Proposition~\ref{prop:galois-s4}).
\end{remark}

\subsection{Monodromy and the generic Galois group}

\begin{proposition}\label{prop:galois-s4}
Viewing $\Pi(\lambda,y)$ as a quartic in $\lambda$ with coefficients in $\QQ(y)$,
its Galois group over $\QQ(y)$ is the full symmetric group:
\[
\mathrm{Gal}\bigl(\Pi(\lambda,y)/\QQ(y)\bigr)\cong S_4.
\]
\end{proposition}

\begin{proof}
Since $\Disc_\lambda \Pi(\lambda,y)=-y(y-1)c(y)$ by Theorem~\ref{thm:disc-factor}, the discriminant is not a square in $\QQ(y)$ (the factors are pairwise coprime irreducibles), so the Galois group is not contained in $A_4$.

To show it is all of $S_4$, specialize at $y=2$: $\Pi(\lambda,2)=\lambda^4-\lambda^3-5\lambda^2+\lambda+6$ is irreducible over $\QQ$ (no rational roots, no factorization into irreducible quadratics).
Reducing modulo $5$: $\Pi(\lambda,2)\equiv \lambda^4+4\lambda^3+4\lambda+1\pmod{5}$, which is irreducible over $\mathbb{F}_5$, giving a Frobenius element of cycle type $(4)$.
Reducing modulo $3$: $\Pi(\lambda,2)\equiv \lambda(\lambda^3+2\lambda^2+\lambda+1)\pmod{3}$, with the cubic factor irreducible over $\mathbb{F}_3$, giving cycle type $(1)(3)$.
Since $\gcd(3\cdot 5,\Disc(\Pi(\lambda,2)))=1$, these Frobenius elements lie in $\mathrm{Gal}(\Pi(\lambda,2)/\QQ)$ by the Chebotarev density theorem.
A subgroup of $S_4$ containing both a $4$-cycle and a $3$-cycle is all of $S_4$.
Since specialization gives a quotient of the generic Galois group,
$\mathrm{Gal}(\Pi(\lambda,y)/\QQ(y))\cong S_4$.
\end{proof}

\begin{remark}[Resolvent cubic]\label{rmk:resolvent}
The \emph{resolvent cubic} of the quartic $\Pi(\lambda,y)$ is the cubic polynomial whose roots are $\lambda_i\lambda_j+\lambda_k\lambda_l$ for the three partitions $\{\{i,j\},\{k,l\}\}$ of $\{1,2,3,4\}$.
For the present spectral polynomial $\Pi(\lambda,y)=\lambda^4+p\lambda^3+q\lambda^2+r\lambda+s$ with $p=-1$, $q=-(2y+1)$, $r=1$, $s=y(y+1)$, the standard formula $R(z)=z^3-qz^2+(pr-4s)z-(p^2 s-4qs+r^2)$ gives
\[
R(z,y)=z^3+(2y+1)z^2-(4y^2+4y+1)z-(8y^3+13y^2+5y+1).
\]
At $y=1$ the resolvent has roots $3,-3,-3$, corresponding to the three pairing sums of the spectral roots $\{2,1,-1,-1\}$.
The Galois group of $R$ over $\QQ(y)$ is $S_3$ (since $\mathrm{Gal}(\Pi/\QQ(y))\cong S_4$ acts on the three roots of $R$ via the quotient $S_4/V_4\cong S_3$).
The splitting field of $R$ corresponds to the $V_4$-quotient of the $S_4$-cover, and its discriminant coincides (up to squares) with the discriminant of $\Pi$.
\end{remark}

\subsection{A divisor-theoretic constraint on fibers}

\begin{corollary}\label{cor:fiber-sum-zero}
Let $y_0\in\CC$ be a non-branch value of $\varphi$ (i.e.\ $y_0\notin\Branch$).
Write $\varphi^{-1}(y_0)=\{P_1,P_2,P_3,P_4\}\subset E(\CC)$.
Then, in the group law on $E$,
\[
P_1+P_2+P_3+P_4=O.
\]
\end{corollary}

\begin{proof}
The rational function $y-y_0$ on $E$ has pole divisor $4O$ by $\operatorname{ord}_O(y)=-4$,
and its zero divisor is $\sum_{i=1}^4 P_i$.
Thus $\operatorname{div}(y-y_0)=\sum_{i=1}^4 P_i-4O$ is principal, hence trivial in $\mathrm{Pic}^0(E)\simeq E$,
giving $\sum_i P_i=O$.
\end{proof}

\subsection{Local square-root branching and the Puiseux amplitude}

\begin{theorem}\label{thm:puiseux-amplitude}
Let $y_\star$ be a simple finite branch value of $\varphi$, and let $P_\star\in E$ be the corresponding ramification point.
Then any local spectral branch $\lambda(y)$ of $\Pi(\lambda,y)=0$ near $y_\star$ has a square-root expansion
\[
\lambda(y)=\lambda_\star \pm \kappa_\star \sqrt{y-y_\star}+O(y-y_\star),
\]
where $(\lambda_\star,y_\star)$ is the branch point.
If $y_\star$ is one of the three nontrivial branch values (i.e.\ $c(y_\star)=0$), then $\lambda_\star$ satisfies
$16\lambda_\star^3-9\lambda_\star^2+1=0$, and the amplitude is
\begin{equation}\label{eq:kappa-formula}
\kappa_\star^{2}=\frac{2(3\lambda_\star^{2}-1)}{3(\lambda_\star^{2}-1)}.
\end{equation}
Moreover, the three values of $\kappa_\star^2$ are the roots of the cubic
\begin{equation}\label{eq:puiseux-cubic}
324u^3-540u^2+333u-74=0,
\end{equation}
whose discriminant is $\Disc=-2^{6}\cdot 3^{9}\cdot 37$.
The branch cubic $c(y)$ has $\Disc(c)=-3^9\cdot 31^2\cdot 37$; both discriminants share the factor $3^9\cdot 37$, reflecting the common arithmetic origin of the two cubics in the ramification structure of $\varphi$.
\end{theorem}

\begin{proof}
The local Puiseux form $y-y_\star=ct^2+O(t^3)$ at a simple ramification point is standard~\cite{Silverman2009},
and inversion gives the asserted square-root structure.

For the explicit amplitude at a nontrivial branch point, the Newton polygon method gives
$\kappa_\star^2=-2\,\partial_y\Pi/\partial^2_{\lambda\lambda}\Pi$ evaluated at $(\lambda_\star,y_\star)$.
With $\partial_y\Pi=-2\lambda^2+2y+1$ and $\partial^2_{\lambda\lambda}\Pi=12\lambda^2-6\lambda-2(2y+1)$,
and using the ramification condition $\xi=(1-2\lambda-3\lambda^2)/(4\lambda)$ together with $16\lambda^3-9\lambda^2+1=0$,
substitution and simplification yields~\eqref{eq:kappa-formula}.

The cubic \eqref{eq:puiseux-cubic} is obtained by eliminating $\lambda_\star$ between \eqref{eq:kappa-formula} and $16\lambda_\star^3-9\lambda_\star^2+1=0$ via resultant computation.
The discriminant of \eqref{eq:puiseux-cubic} is $-2^{6}\cdot 3^{9}\cdot 37$, computed directly from the standard cubic discriminant formula.
\end{proof}

\subsection{Modularity}

\begin{theorem}[Modularity at level $37$]\label{thm:modularity}
Let $E/\QQ$ be the elliptic curve \eqref{eq:elliptic-E}.
Then there exists a weight-$2$ newform $f\in S_2(\Gamma_0(37))$ such that $L(E,s)=L(f,s)$.
In particular, for each prime $p\neq 37$, the number of $\mathbb{F}_p$-points on $E$ satisfies
\begin{equation}\label{eq:frobenius-trace}
\#E(\mathbb{F}_p)=p+1-a_p,
\end{equation}
where $a_p$ is the $p$-th Fourier coefficient of $f$.
\end{theorem}

\begin{proof}
By the modularity theorem (Wiles~\cite{WilesTaylor1995}, completed by Breuil--Conrad--Diamond--Taylor), every elliptic curve over $\QQ$ is modular.
The level of the associated newform equals the conductor, which is $37$ by Theorem~\ref{thm:ellipticization}.
The point-counting formula \eqref{eq:frobenius-trace} is the Hasse--Weil relation, with $a_p=p+1-\#E(\mathbb{F}_p)$ matching the Fourier coefficients of $f$ by the Eichler--Shimura theorem~\cite{Silverman2009}.
\end{proof}

\begin{remark}[Finite-field spectral interpretation]\label{rmk:finite-field}
Since $C$ and $E$ are isomorphic over $\QQ$ (Theorem~\ref{thm:ellipticization}), the $\mathbb{F}_p$-points of the spectral curve $\Pi(\lambda,y)\equiv 0\pmod{p}$ are in bijection with $E(\mathbb{F}_p)$ for all primes $p\neq 37$.
Thus the spectral polynomial $\Pi$ ``knows'' the Frobenius traces $a_p$ of the modular form, via
$\#\{(\lambda,y)\in\mathbb{F}_p^2:\Pi(\lambda,y)\equiv 0\}=p-a_p$.
\end{remark}

\subsection{Arithmetic of the branch cubic field}

\begin{proposition}\label{prop:cubic-disc}
Let $c(y)=256y^3+411y^2+165y+32$ be the branch cubic from \eqref{eq:cubic}, and let $K=\QQ(\yLY)$ be the cubic number field generated by its real root.
\begin{enumerate}
\item[(i)] The polynomial discriminant is $\Disc(c)=-3^{9}\cdot 31^{2}\cdot 37$.
\item[(ii)] The field discriminant is $\Disc(K/\QQ)=-999=-3^3\cdot 37$.
\item[(iii)] The index satisfies $[\mathcal{O}_K:\ZZ[\yLY]]=3^3\cdot 31=837$.
\item[(iv)] The unique quadratic subfield of the Galois closure of $K/\QQ$ is $\QQ(\sqrt{-111})=\QQ(\sqrt{-3\cdot 37})$.
\end{enumerate}
\end{proposition}

\begin{proof}
(i) The standard cubic discriminant formula applied to $(a,b,c,d)=(256,411,165,32)$ gives
$\Disc(c)=b^2c^2-4ac^3-4b^3d-27a^2d^2+18abcd=-3^9\cdot 31^2\cdot 37$.

(ii)--(iii) The relation between polynomial and field discriminants is $\Disc(c)=[\mathcal{O}_K:\ZZ[\alpha]]^2\cdot\Disc(K/\QQ)$ for a root $\alpha$ of $c$~\cite{Silverman2009}.
The field discriminant and index can be determined by computing the $p$-adic valuations: at $p=3$, the polynomial discriminant has $v_3=9$, and the maximal order computation gives $v_3(\Disc(K))=3$, hence $v_3([\mathcal{O}_K:\ZZ[\yLY]])=3$; at $p=31$, $v_{31}(\Disc(c))=2$ and $v_{31}(\Disc(K))=0$, giving $v_{31}([\mathcal{O}_K:\ZZ[\yLY]])=1$.
All other primes divide $\Disc(c)$ to at most the first power, so $[\mathcal{O}_K:\ZZ[\yLY]]=3^3\cdot 31=837$ and $\Disc(K)=-3^3\cdot 37=-999$.

(iv) The Galois closure of $K/\QQ$ has Galois group $S_3$ (since $c$ is irreducible with non-square discriminant).
The unique quadratic subfield corresponds to the sign character $S_3\to\{\pm 1\}$ and is $\QQ(\sqrt{\Disc(K)})=\QQ(\sqrt{-999})=\QQ(\sqrt{-111})$.
\end{proof}

\begin{remark}
The appearance of $37$ in both the elliptic-curve conductor ($\Delta(E)=37$) and the cubic field discriminant ($\Disc(K)\propto 37$) reflects a structural relation: the branch cubic $c(y)$ is determined by the ramification of the degree-$4$ map $\varphi:E\to\mathbb{P}^1$, and the arithmetic of $E$ (particularly its bad reduction at $p=37$) propagates into the discriminant of $c$.
This ``$37$-coincidence'' is not accidental but a consequence of the shared algebraic origin in the spectral curve.
\end{remark}

